\subsection{Preguntas e hipótesis experimentales}

La pregunta central se dividirá en comparaciones medibles sobre especialización,
transferencia, imitación, generalización y memoria recurrente. Cada pregunta se
asociará antes del análisis con métricas y evidencia esperada.

\subsection{Diseño global de la campaña}

Se presentará el flujo completo entre verificación, generación de datos,
sintonización, selección de demostraciones, entrenamiento y evaluación. También
se explicará cómo se evita utilizar prueba para tomar decisiones.

\subsection{Diseño del conjunto de datos clásico}

Se justificarán familias, geometrías, perturbaciones, número de episodios,
semilla y particiones de entrenamiento, validación y prueba. No bastará con
describir los 150 escenarios: deberá explicarse por qué son adecuados.

\subsection{Diseño del conjunto de datos de imitación}

Se mostrará cómo cada escenario clásico se transforma en una demostración de
fuerza deseada segura y trazable. Se explicará qué distribución de
comportamientos aprende realmente la red.

\subsection{Configuración y comparación del entrenamiento}

Se fijarán hiperparámetros, semillas, dispositivo y criterios de selección para
comparar MLP, GRU y LSTM bajo condiciones equivalentes. Se documentará la
experimentación utilizada para tomar estas decisiones.

\subsection{Evaluación en familias conocidas y transferencia clásica}

Se definirá la comparación dentro de la distribución de entrenamiento (ID) y la
matriz cruzada de controladores PD. Este nivel establecerá el rendimiento
especializado y la referencia que pretende condensar la red.

\subsection{Evaluación fuera de la distribución de entrenamiento}

\subsubsection{Variaciones y composiciones de familias conocidas}

Se evaluarán condiciones que combinan o modifican capacidades presentes en el
entrenamiento. Cada escenario deberá indicar qué componentes son conocidos y
qué novedad introduce.

\subsubsection{Trayectorias completamente nuevas}

Se evaluarán familias geométricas nuevas, separándolas de las composiciones para
no mezclar niveles distintos de generalización.

\subsection{Métricas, criterios de éxito y análisis}

Se definirán métricas de seguimiento, actuación, saturación, degradación,
protecciones, éxito de misión y seguridad. Cada métrica incluirá unidad,
interpretación y limitaciones.

\subsection{Reproducibilidad y validez de la evidencia}

Se establecerán los comandos, configuraciones, hashes y metadatos necesarios
para aceptar un resultado. También se explicarán los criterios para excluir
ejecuciones incompletas o no comparables.
